\problem{Archive with Tombstone Compaction}

An online news platform stores articles in an array. Each array entry is either a live article or a tombstone. A tombstone means that the article has been deleted, but its array position has not yet been physically removed.

The system also maintains a mapping from article IDs to array positions. Therefore, checking whether an ID corresponds to a live article and marking it as a tombstone can both be done in \(O(1)\) time. During compaction, the system also updates this mapping. Updating the mapping for a copied live article is included in the copying cost of that live article.

Users can perform the following two external operations:
\begin{itemize}
    \item \textsc{Add}\((\textit{article})\): append the new article as a live article at the end of the array, and record the mapping from its ID to its array position.
    \item \textsc{Remove}\((\textit{id})\): if the article with ID \(\textit{id}\) is still live, mark its entry as a tombstone. If this ID does not exist or is no longer live, perform only one lookup and return.
\end{itemize}

The system also has the following internal cleanup procedure:
\begin{itemize}
    \item \textsc{Compact}\(()\): scan the entire array, copy all live articles into a new array, discard all tombstones, and update the ID-to-position mapping.
\end{itemize}
\textsc{Compact}\(()\) is not directly called by users. It is automatically triggered by the system after a \textsc{Remove} operation.

System rule: after each \textsc{Remove}, if the number of tombstones is strictly greater than the number of live articles, the system automatically executes \textsc{Compact}\(()\).

Assume that appending an article costs \(1\), marking a tombstone costs \(1\), and scanning or copying one array entry costs \(1\).

Answer the following questions:
\begin{enumerate}[(a)]
    \item Explain why a single \textsc{Remove} operation can have worst-case cost \(\Theta(n)\), where \(n\) is the array length immediately before the \textsc{Compact} operation triggered by this \textsc{Remove}.
    \item Prove that the total cost of any \(m\) external operations is \(O(m)\).
    \item Give a potential function based on the number of tombstones and prove an amortized \(O(1)\) bound.
\end{enumerate}

\solution{
\textbf{(a)}
A \textsc{Remove} operation may trigger \textsc{Compact}. If the array length is
\(n\) immediately before this compaction, then \textsc{Compact} scans the whole
array, which already costs \(\Theta(n)\). It may also copy many live articles.
Hence a single \textsc{Remove} operation can have worst-case cost
\(\Theta(n)\).

\textbf{(b)}
Let \(D\) be the number of tombstones and \(L\) be the number of live articles
immediately after a \textsc{Remove} operation and before a possible compaction.
Compaction is triggered only when
\[
    D>L.
\]
The cost of this compaction is at most the cost of scanning all entries plus
the cost of copying all live entries:
\[
    (D+L)+L = D+2L.
\]
Since \(D>L\), this is less than \(3D\).

Charge a constant number of credits to every successful \textsc{Remove} that
creates a tombstone. One part pays for the constant cost of this
\textsc{Remove}; the remaining credits stay with the tombstone. Each tombstone
is discarded at the next compaction, and the above inequality shows that the
tombstones present at that time can pay for the whole compaction cost. Each
\textsc{Add} costs \(O(1)\), and a \textsc{Remove} of a nonexistent or already
deleted ID also costs \(O(1)\). Therefore the total cost of any \(m\) external
operations is \(O(m)\).

\textbf{(c)}
Let
\[
    \Phi = 3D,
\]
where \(D\) is the current number of tombstones in the array. Initially
\(\Phi=0\), and \(\Phi\ge 0\) always.

For \textsc{Add}, the actual cost is \(O(1)\), and \(D\) does not change. Thus
the amortized cost is \(O(1)\).

For an unsuccessful \textsc{Remove}, the operation only does a lookup and
returns. Again \(D\) does not change, so the amortized cost is \(O(1)\).

For a successful \textsc{Remove} that does not trigger compaction, the actual
cost is \(O(1)\), and \(D\) increases by \(1\). Hence
\[
    \hat c = O(1)+3=O(1).
\]

Now consider a successful \textsc{Remove} that triggers compaction. Let \(D'\)
and \(L'\) be the numbers of tombstones and live articles after the tombstone
is created but before compaction. Then \(D'>L'\). The compaction cost is at
most
\[
    D'+2L' < 3D'.
\]
Before this \textsc{Remove}, the number of tombstones was \(D'-1\), so the
potential was \(3(D'-1)\). After compaction, the number of tombstones is \(0\),
so the potential becomes \(0\). Therefore the amortized cost is at most
\[
    O(1) + (D'+2L') - 3(D'-1)
    < O(1)+3.
\]
This is \(O(1)\). Hence every external operation has amortized cost \(O(1)\).
}
\newpage
